\section{Conclusion}
\label{sec:conclusion}

Sparse Concept Anchoring inverts traditional interpretability: rather than discovering concept locations post-hoc, it establishes predictable geometric structure during training through minimal supervision. By fixing concept directions and regularizing representations, SCA enables direct interventions on model behavior.

Experiments on color autoencoders show that both inference-time suppression and parameter-level ablation achieve selective concept modification with minimal collateral impact. Sparse, noisy labels ($<0.1\%$ of data) shaped latent geometry into actionable subspaces, demonstrating that anchoring few concepts suffices for controllable representations while preserving flexibility.

Although current experiments use low-dimensional settings, geometric regularization with sparse supervision provides a path toward scalable concept-level control. For safety-critical applications requiring auditable, modifiable systems, proactively shaping representations may prove more tractable than reactive analysis of emergent structure.
